% app_extra_B.tex -- Appendix C, part B (author B): supplements to Section 5, not used by the main theorem.

\subsection{Lazy monotonicity and the eager variant}\label{sec:extra-lazy}

This subsection contains the proofs of \cref{prop:lazy-monotone} and
\cref{cor:eager}. The notation is that of \cref{sec:game}; the label \WON{} refers to a fixed
$\Rmax$, and $\rk$ is the rank of \cref{lem:rank}. Neither statement is used in the proof of
\cref{thm:main}.

\begin{lemma}[$\LOST$ is antitone in $F$]\label{lem:lost-antitone}
Let $F\subseteq F'$ be sets of cells disjoint from $B$. If $\LOST(B,F',t)$ holds, then $\LOST(B,F,t)$ holds.
\end{lemma}

\begin{proof}
Let a pattern for $(B,F')$ fit $(B,F',t)$. The same lines, cells and stone sets form a pattern for $(B,F)$:
the conditions (E2.1), (E2.2), (E3.1), (E3.2) of \cref{def:lost} do not involve the revealed set,
and the conditions (E2.3), (E3.3), which are of the form ``$e\notin F'$'' or ``$F'\cap a\subseteq S_a$'', remain
true when $F'$ is replaced by a subset. The union of $F$ and the stone sets is contained in the union of
$F'$ and the stone sets, so it has at most $t+1$ cells.
\end{proof}

\begin{proof}[Proof of \cref{prop:lazy-monotone}]
The set $F$ is $B$-admissible, because $F\subseteq F'$ and $F'$ is $B$-admissible
(\cref{lem:adm-closed}\,(a)). Hence none of the four states that occur in (a) and (b) is \VOID. We prove (a)
and (b) simultaneously by induction on the rank of the state with the smaller set $F$, that is, of
$\Pst(B,F,t)$ in (a) and of $\Qst(B,F,t)$ in (b).

(a) Let $\Pst(B,F,t)$ be \WON. If $\LEAF(B,F,t)$ holds, then $\LEAF(B,F',t)$ holds by the monotonicity of
$\LEAF$ in the revealed set (\cref{prop:leaf-monotone}\,(a), which applies because $F'$ is
$B$-admissible and $|F'|\le t$), and $\Pst(B,F',t)$ satisfies (P1). Otherwise $\Pst(B,F,t)$ satisfies (P2):
$t<\Rmax$, $\LOST(B,F,t)$ is false, and every child of $\Pst(B,F,t)$ is \WON{} or \VOID. We verify (P2) for
$\Pst(B,F',t)$. First, $\LOST(B,F',t)$ is false by \cref{lem:lost-antitone}.

\emph{Child \textup{(a)}.} The child $\Qst(B,F,t+1)$ of $\Pst(B,F,t)$ is not \VOID, hence \WON. It has
smaller rank than $\Pst(B,F,t)$, and $|F'|\le t+1$. By induction, part (b), $\Qst(B,F',t+1)$ is \WON.

\emph{Children \textup{(c)}.} Let $l\in\live{B}$ and $e\in l\setminus F'$ with
$|F'\cup(l\setminus\{e\})|\le t+1$ define a child \textup{(c)} of $\Pst(B,F',t)$, namely
$\Pst(B+e,F_1',t+1)$ with $F_1'=F'\cup(l\setminus\{e\})$, and suppose that this child is not \VOID, that is,
$F_1'$ is $(B+e)$-admissible. Since $e\notin F$ and
$|F\cup(l\setminus\{e\})|\le|F'\cup(l\setminus\{e\})|\le t+1$, the same data define a child \textup{(c)}
of $\Pst(B,F,t)$, namely $\Pst(B+e,F_1,t+1)$ with $F_1=F\cup(l\setminus\{e\})\subseteq F_1'$. This child is
not \VOID{} (\cref{lem:adm-closed}\,(a)), hence it is \WON, and it has smaller rank than $\Pst(B,F,t)$. As
$F_1\subseteq F_1'$, $F_1'$ is $(B+e)$-admissible and $|F_1'|\le t+1$, the state $\Pst(B+e,F_1',t+1)$ is \WON{}
by induction, part (a).

\emph{Children \textup{(d)}.} Let $a,b,c,Q_a,Q_b$ define a child \textup{(d)} of $\Pst(B,F',t)$ that is not
\VOID. From $c\notin F'$, $F'\cap a\subseteq Q_a$, $F'\cap b\subseteq Q_b$ and $|F'\cup Q_a\cup Q_b|\le t+1$ we
get the same four conditions for $F$. So the same data define a child \textup{(d)} of $\Pst(B,F,t)$, whose
revealed set $F\cup Q_a\cup Q_b$ is contained in $F'\cup Q_a\cup Q_b$. The argument is completed as for
\textup{(c)}.

Thus every child of $\Pst(B,F',t)$ is \WON{} or \VOID, and $\Pst(B,F',t)$ satisfies (P2).

(b) Let $\Qst(B,F,t)$ be \WON, and let $m\notin B\cup F$ be a table move for it.

\emph{Case $m\in F'$.} Then $m\in F'\setminus F$, so $|F|<|F'|\le t$, and the child \textup{(ii)} of
$\Qst(B,F,t)$ via $m$, the state $\Qst(B,F+m,t)$, exists. It is not \VOID, because $F+m\subseteq F'$; hence
it is \WON{} by requirement (ii) of \cref{def:won}. Its rank is smaller than that of $\Qst(B,F,t)$,
and $F+m\subseteq F'$. By induction, part (b), $\Qst(B,F',t)$ is \WON.

\emph{Case $m\notin F'$.} We show that $m$ is a table move for $\Qst(B,F',t)$; note $m\notin B\cup F'$.
Requirement (i): the state $\Pst(B+m,F,t)$ is \WON{} and has smaller rank than $\Qst(B,F,t)$; the set $F'$ is
$(B+m)$-admissible because $m\notin F'$ (\cref{lem:adm-closed}\,(b)), and $|F'|\le t$. By induction, part
(a), $\Pst(B+m,F',t)$ is \WON. Requirement (ii): suppose that $|F'|<t$ and that $\Qst(B,F'+m,t)$ is not
\VOID, that is, $F'+m$ is $B$-admissible. Then $|F|<t$, so the child \textup{(ii)} of $\Qst(B,F,t)$ via $m$,
the state $\Qst(B,F+m,t)$, exists; it is not \VOID, because $F+m\subseteq F'+m$; hence it is \WON, and its
rank is smaller than that of $\Qst(B,F,t)$. As $F+m\subseteq F'+m$ and $|F'+m|\le t$, the state
$\Qst(B,F'+m,t)$ is \WON{} by induction, part (b).
\end{proof}

\begin{proof}[Proof of \cref{cor:eager}]
We show by induction on the rank that every state that is \WON{} is $\WONp$. Let $\Pst(B,F,t)$ be \WON. If
it satisfies (P1), it is $\WONp$, since (P1) is the same for both labels. Otherwise it satisfies (P2); its
children \textup{(a)}, \textup{(c)}, \textup{(d)} are \WON{} or \VOID, hence $\WONp$ or \VOID{} by induction.
It remains to consider a child \textup{(b)}, that is, $\Qst(B,F+h,t+1)$ with $h\in H\setminus(B\cup F)$. If
$F+h$ is not $B$-admissible, this state is \VOID. Otherwise observe that the child \textup{(a)},
$\Qst(B,F,t+1)$, is \WON, and that $F\subseteq F+h$, where $F+h$ is $B$-admissible and $|F+h|\le t+1$. By
\cref{prop:lazy-monotone}\,(b), $\Qst(B,F+h,t+1)$ is \WON. Its rank is smaller than that of
$\Pst(B,F,t)$, so it is $\WONp$ by induction. Finally let $\Qst(B,F,t)$ be \WON, with table move $m$. Its
child \textup{(i)} via $m$ is \WON, hence $\WONp$ by induction, and its child \textup{(ii)} via $m$, if it
exists, is \WON{} or \VOID, hence $\WONp$ or \VOID. So $\Qst(B,F,t)$ is $\WONp$.
\end{proof}

\subsection{The patterns \texorpdfstring{$\Etwo$}{E2} and \texorpdfstring{$\Ethree$}{E3} on the
\texorpdfstring{$6\times6\times6$}{6x6x6} board}\label{sec:extra-e23}

\cref{prop:lost-stones} bounds the number of distinct stones of a pattern from below by $11$ for
$\Etwo$ and by $12$ for $\Ethree$, using only that two lines share at most one cell. On the
$6\times6\times6$ board the first bound is sharp and the second is not. Nothing in the paper depends on
this.

For the proof we regard cells as points of $\mathbb{R}^3$. Let
$D=\{-1,0,1\}^3\setminus\{(0,0,0)\}$. Recall that a line is a set $\{(f_1(t),f_2(t),f_3(t)):0\le t\le 5\}$,
where each $f_i$ is a constant, $t$, or $5-t$, and not all three are constant. Such a line lies on a real
line with direction vector $d\in D$, where $d_i$ is $0$, $1$ or $-1$ according as $f_i$ is constant, $t$ or
$5-t$; the vector $d$ is determined by the line up to sign. If $d_i=d_k\ne 0$, then all cells
$q=(q_1,q_2,q_3)$ of the line satisfy $q_i=q_k$, and if $d_i=-d_k\ne 0$, they satisfy $q_i+q_k=5$. A real line
with a direction vector in $D$ contains at most six cells of the board, so distinct lines of the board lie
on distinct real lines.

\begin{proposition}[Patterns on the $6\times6\times6$ board]\label{prop:e23-min}
On the $6\times6\times6$ board:
\begin{enumerate}[label=\textup{(\alph*)},leftmargin=2.4em]
\item any four lines that meet pairwise pass through one common cell;
\item every $\Ethree$ pattern has at least $13$ distinct stones, and there are $\Ethree$ patterns with
  exactly $13$;
\item there are $\Etwo$ patterns with exactly $11$ distinct stones.
\end{enumerate}
\end{proposition}

\begin{proof}
(a) Let $a_1,a_2,a_3,a_4$ be lines that meet pairwise, and suppose that they do not pass through a common
cell. Distinct lines that meet lie on distinct real lines through a common point, so the direction vectors
$d^1,\dots,d^4\in D$ of the four lines are pairwise non-proportional.

We first show that the four lines lie in a plane. If no three of them pass through a common cell, then $a_1$
and $a_2$ span a plane, and each of $a_3$, $a_4$ meets $a_1$ and $a_2$ in two different cells and therefore
lies in this plane. If three of them, say $a_1,a_2,a_3$, pass through a common cell $p$, then $a_4$ does not
contain $p$ and meets $a_1,a_2,a_3$ in three different cells, which are collinear; so $a_1,a_2,a_3$ lie in the
plane spanned by $p$ and $a_4$. Hence $d^1,\dots,d^4$ lie in a two-dimensional subspace $W$ of $\mathbb{R}^3$.

We claim that $W$ contains a vector with exactly one non-zero coordinate. Suppose not. The \emph{support}
of a vector is the set of its non-zero coordinates. Let $u,v\in D\cap W$ be non-proportional vectors with
the same support. In every coordinate of the support $u_i=\pm v_i$, so exactly one of $(u+v)_i$, $(u-v)_i$
is non-zero, and it equals $\pm2$. Thus $u+v$ and $u-v$ are non-zero vectors of $W$ whose supports partition
the common support of $u$ and $v$, which has at most three elements; one of them has a support of size one, a
contradiction. So the supports of $d^1,\dots,d^4$ are four different sets of size two or three, that is, they
are $\{1,2\}$, $\{1,3\}$, $\{2,3\}$ and $\{1,2,3\}$. Let $u,v,w$ be the three of these vectors with supports
$\{1,2\}$, $\{1,3\}$, $\{1,2,3\}$. Then $w=\alpha u+\beta v$ with real $\alpha,\beta$; the second coordinate
gives $\alpha=\pm1$, the third gives $\beta=\pm1$, and then the first coordinate
$w_1=\alpha u_1+\beta v_1$ is even, a contradiction. This proves the claim.

So $W$ contains a unit vector, say the one of coordinate $k$; let $i,j$ be the other two coordinates. The
vectors of $W$ are then the vectors $z$ for which $(z_i,z_j)$ is a multiple of one fixed non-zero pair,
which can be taken from one of the vectors $d^1,\dots,d^4$ and thus, up to sign, is $(1,0)$, $(0,1)$,
$(1,1)$ or $(1,-1)$. In the first two cases $W$ is a coordinate plane, say $W=\{z: z_j=0\}$, and the vectors
of $D\cap W$ are, up to sign, the two unit vectors of the coordinates $i$ and $k$ and two vectors $u,u'$
with $u_i=u_k\ne0$ and $u'_i=-u'_k\ne0$. In the last two cases $W=\{z: z_i=\varepsilon z_j\}$ with
$\varepsilon=\pm1$, and the vectors of $D\cap W$ are, up to sign, the unit vector of coordinate $k$, one
vector with $z_k=0$, and two vectors $u,u'$ with $u_i=u_k\ne0$ and $u'_i=-u'_k\ne0$. In every case $D\cap W$
consists of exactly four vectors up to sign, so two of the four lines have the direction vectors $u$ and
$u'$. The cells $q$ of the first of them satisfy $q_i=q_k$, and those of the second satisfy $q_i+q_k=5$. A
common cell would satisfy $2q_i=5$, which is impossible. So these two lines do not meet, contrary to the
assumption.

(b) Suppose an $\Ethree$ pattern has only $12$ distinct stones. By \cref{prop:lost-stones}\,(b) the
sets $S_a\cap S_{b'}$, $S_{a'}\cap S_b$ and $S_b\cap S_{b'}$ are then non-empty, so $a$ meets $b'$, $a'$ meets
$b$, and $b$ meets $b'$. By (E3.1), $a\cap a'=\{x\}$, $a\cap b=\{c\}$ and $a'\cap b'=\{c'\}$. So the four lines
meet pairwise, and by (a) they pass through a common cell. This cell is $x$, the only common cell of $a$ and
$a'$; but $x\notin b$ by (E3.1), a contradiction. An $\Ethree$ pattern with $13$ distinct stones is the
following one, in the plane $z=0$; we write $(i,j)$ for the cell $(i,j,0)$. Lines:
$a=\{(i,0)\}$, $a'=\{(0,j)\}$, $b=\{(i,5-i)\}$, $b'=\{(i,1)\}$; cells: $x=(0,0)$, $c=(5,0)$, $c'=(0,1)$;
stone sets:
\begin{align*}
S_a&=\{(0,0),(1,0),(2,0),(3,0)\}, & S_{a'}&=\{(0,0),(0,2),(0,3),(0,5)\},\\
S_b&=\{(0,5),(1,4),(2,3),(4,1)\}, & S_{b'}&=\{(1,1),(2,1),(3,1),(4,1)\}.
\end{align*}
The conditions (E3.1) and (E3.2) are checked directly, and (E3.3) holds for $F=\emptyset$. The stone sets
meet only in $S_a\cap S_{a'}=\{(0,0)\}$, $S_{a'}\cap S_b=\{(0,5)\}$ and $S_b\cap S_{b'}=\{(4,1)\}$, so the
pattern has $16-3=13$ distinct stones.

(c) With the same notation let $l=\{(i,0)\}$, $a=\{(0,j)\}$, $b=\{(i,5-i)\}$, $x=(0,0)$, $e=(1,0)$,
$c=(0,5)$, $S_a=\{(0,0),(0,1),(0,2),(0,3)\}$ and $S_b=\{(5,0),(4,1),(3,2),(2,3)\}$. The conditions (E2.1)
and (E2.2) are checked directly, and (E2.3) holds for $F=\emptyset$. Here
$S_l=l\setminus\{e\}$ meets $S_b$ in $(5,0)$, so the pattern has $12-1=11$ distinct stones by
\cref{prop:lost-stones}\,(a).
\end{proof}

The exact distribution of the number of stones was determined by computer.

\begin{fact}[Stones of the patterns on the empty board]\label{fact:e23}
For $B=F=\emptyset$, when all $148$ lines are live, there are $5{,}152{,}320$ $\Etwo$ patterns, of which
$562{,}176$ have $11$ and $4{,}590{,}144$ have $12$ distinct stones. There are $353{,}422{,}080$ $\Ethree$
patterns, of which $10{,}948{,}608$ have $13$, $67{,}627{,}008$ have $14$, and $274{,}846{,}464$ have $15$ distinct
stones; none has $12$. Here a pattern is the tuple of its lines, cells and stone sets as listed in
\cref{def:lost}, so that an $\Ethree$ pattern and the pattern obtained from it by exchanging
$(a,b,c,S_a,S_b)$ with $(a',b',c',S_{a'},S_{b'})$ are counted separately.

\established the program \fn{e23min.cpp} in the directory \fn{logs/proof/} of the project
repository (output \fn{e23min.log}; a few seconds of running time; see \cref{ssec:artifact}). The program is self-contained and constructs the
lines itself; it enumerates all patterns literally from \cref{def:lost} and counts the cells in the
union of the stone sets.
\end{fact}

Every pattern for an arbitrary pair $(B,F)$ is also a pattern for $(\emptyset,\emptyset)$, because a
non-empty $B$ only removes live lines and a non-empty $F$ only restricts the stone sets. So \cref{fact:e23} covers all patterns; it agrees with \cref{prop:lost-stones} and with
\cref{prop:e23-min}.

\smallskip
The last example of this subsection belongs to \cref{rem:conservative}(ii): it exhibits a position in
which an $\Ethree$ pattern fits although \Breaker{} has a reply that restores admissibility.

\begin{example}[A fitting pattern that does no harm]\label{ex:conservative}
In the plane $z=0$ consider the lines
\[ a=\{(i,0,0)\},\qquad a'=\{(0,j,0)\},\qquad b=\{(1,j,0)\},\qquad b'=\{(i,5-i,0)\}, \]
where $i$ and $j$ run from $0$ to $5$, the cell $x=(0,0,0)$, and the set $X'$ of the $14$ cells $(i,j,0)$
with
\begin{align*}
(i,j)\in\{&(0,0),(2,0),(3,0),(4,0),\ (0,1),(0,2),(0,3),\\
          &(1,1),(1,2),(1,3),(1,4),\ (2,3),(3,2),(4,1)\}.
\end{align*}
Let $X=X'\setminus\{x\}$, let $B$ be any set of $13$ cells with $z\ne 0$, and let $F\subseteq X$ be arbitrary;
then $\Pst(B,F,13)$ is a P-state. A line that is not contained in the plane $z=0$ has at most one cell in
this plane. The $14$ lines in the plane are live. The six lines parallel to $a$ (given by $j=0,\dots,5$) have
$4,3,3,3,1,0$ cells in $X'$, the six lines parallel to $a'$ (given by $i=0,\dots,5$) have $4,4,2,2,2,0$, the
diagonal $\{(i,i,0)\}$ has two, and $b'$ has four. So $a$, $a'$, $b$, $b'$ have four cells in $X'$ each, and
every other line has at most three. In $X$ the lines $a$ and $a'$ have three cells, and $b$, $b'$ have four
cells and meet in $(1,4,0)\in X$. So $X$ is $B$-admissible, and $X'=X+x$ has no threat. The lines with four
cells in $X'$ meet as follows: $a\cap a'=\{x\}$ and $b\cap b'=\{(1,4,0)\}$ are in $X'$,
$a'\cap b=\emptyset$, and
\[ a\cap b=\{(1,0,0)\},\qquad a\cap b'=\{(5,0,0)\},\qquad a'\cap b'=\{(0,5,0)\} \]
are empty cells. So $X'$ has exactly the danger pairs $(a,b)$, $(a,b')$ and $(a',b')$. By
\cref{lem:danger}, applied to $(a,b)$ and $(a',b')$, an $\Ethree$ pattern fits $(B,F,13)$, and indeed the
replies $(1,0,0)$ and $(0,5,0)$ leave the danger pairs $(a',b')$ and $(a,b)$, respectively. But after the
reply $(5,0,0)$ only $a'$ and $b$ are live lines with four cells in $X'$, and they are disjoint: $X'$ is
$(B+(5,0,0))$-admissible. (Rule \ref{S:danger} of \cref{def:sigma} would choose $(0,5,0)$, the common
cell with the smallest id. The choice of the danger pair can matter, but only at states with $\LOST$, which a
play guided by a valid certificate never reaches.)
\end{example}
